\chapter{ORTHOGONAL FUNCTIONS WHICH SATISFY FOUR BOUNDARY CONDITIONS}
\label{app:V}

\section{Introduction}
\label{sec:132}

In the solution of hydrodynamic problems it is often convenient to be
able to expand the prevailing velocity fields in a complete set of orthogonal functions. However, the usual expansions in terms of trigonometric or Bessel functions are not useful (if not impossible) in problems
involving viscous flow since, in these problems, the boundary conditions
require that on a rigid wall, for example, not only the normal component
of the velocity, but also its normal derivative, vanishes (the latter
condition ensuring the vanishing of the parallel component of the
velocity). In these problems, expansions must therefore be sought in
terms of orthogonal functions which together with their first derivatives
vanish at the ends of the chosen interval. In this appendix, we shall
give an account of the functions which have been constructed for these
applications, though their completeness for the purposes of expansion
requires a more careful analysis than has been given.

\section{Functions suitable for problems with plane boundaries}
\label{sec:133}

Consider the characteristic value problem defined by the equation
\begin{equation}
\label{eq:A5-1}
\frac{d^4 y}{dx^4} = \alpha y
\end{equation}
and the boundary conditions
\begin{equation}
\label{eq:A5-2}
y = 0 \quad \text{and} \quad dy/dx = 0 \quad \text{for} \quad x = \pm\tfrac{1}{2}.
\end{equation}

Let $\alpha_n$ denote a characteristic value and let the proper solution belonging to it be distinguished by the subscript~$n$. Then multiplying the
equation governing $y_n$ by $y_m$ (belonging to $\alpha_m$) and integrating over the
range of $x$, we obtain
\begin{equation}
\label{eq:A5-3}
\alpha_k \int_{-\frac{1}{2}}^{+\frac{1}{2}} y_n \, y_m \, dx = \int_{-\frac{1}{2}}^{+\frac{1}{2}} y_m \frac{d^4 y_n}{dx^4} \, dx.
\end{equation}

Transforming the integral on the right-hand side by two successive
integrations by parts, we are left with
\begin{equation}
\label{eq:A5-4}
\alpha_k \int_{-\frac{1}{2}}^{+\frac{1}{2}} y_n \, y_m \, dx = \int_{-\frac{1}{2}}^{+\frac{1}{2}} \frac{d^2 y_n}{dx^2} \frac{d^2 y_m}{dx^2} \, dx,
\end{equation}
the integrated parts vanishing (both times) on account of the boundary
conditions of $y_n$. From the symmetry of the integral on the right-hand
side of equation~(\ref{eq:A5-4}) in $n$ and $m$, it follows that
\begin{equation}
\label{eq:A5-5}
(\alpha_n - \alpha_m) \int_{-\frac{1}{2}}^{+\frac{1}{2}} y_n \, y_m \, dx = 0.
\end{equation}

Hence
\begin{equation}
\label{eq:A5-6}
\int_{-\frac{1}{2}}^{+\frac{1}{2}} y_n \, y_m \, dx = 0 \quad \text{if} \quad m \neq n.
\end{equation}

The solutions $y_n$, therefore, form an orthogonal set; and we shall suppose
that they are complete for the purposes of expanding an arbitrary
continuous function in the interval $(-\frac{1}{2}, +\frac{1}{2})$ which together with its
first derivative vanishes at the ends.

The proper solutions of equations~(\ref{eq:A5-1}) and~(\ref{eq:A5-2}) are readily found. They
fall into two non-combining groups of even and odd solutions with
cosine-like and sine-like behaviours. As standard forms of these solutions we shall take
\begin{equation}
\label{eq:A5-7}
C_n(x) = \frac{\cosh \lambda_n x}{\cosh \frac{1}{2}\lambda_n} - \frac{\cos \lambda_n x}{\cos \frac{1}{2}\lambda_n}
\end{equation}
and
\begin{equation}
\label{eq:A5-8}
S_n(x) = \frac{\sinh \mu_n x}{\sinh \frac{1}{2}\mu_n} - \frac{\sin \mu_n x}{\sin \frac{1}{2}\mu_n}.
\end{equation}
Defined in this manner, the functions clearly vanish at $x = \pm \frac{1}{2}$; that
their derivatives also vanish at these points requires that $\lambda_n$ and $\mu_n$
are the roots of the characteristic equations
\begin{equation}
\label{eq:A5-9}
\tanh \tfrac{1}{2}\lambda + \tan \tfrac{1}{2}\lambda = 0
\end{equation}
and
\begin{equation}
\label{eq:A5-10}
\coth \tfrac{1}{2}\mu - \cot \tfrac{1}{2}\mu = 0,
\end{equation}
respectively.

It can be readily verified that the functions $C_n(x)$ and $S_n(x)$ are
already normalized so that
\begin{equation}
\label{eq:A5-11}
\int_{-\frac{1}{2}}^{+\frac{1}{2}} C_n(x) C_m(x) \, dx = \int_{-\frac{1}{2}}^{+\frac{1}{2}} S_n(x) S_m(x) \, dx = \delta_{nm}.
\end{equation}
Also, it is evident that every function $C_n$ is orthogonal to every function $S_m$; thus
\begin{equation}
\label{eq:A5-12}
\int_{-\frac{1}{2}}^{+\frac{1}{2}} C_n(x) S_m(x) \, dx = 0.
\end{equation}

In Table~LXVIII the first few roots of equations~(\ref{eq:A5-9}) and~(\ref{eq:A5-10}) and
some related constants are listed.

\begin{table}[htbp]
\centering
\caption{The characteristic roots $\lambda_n$ and $\mu_n$ and related constants}
\label{tab:LXVIII}
\begin{tabular}{lccc}
\hline
 & $n=1$ & $n=2$ & $n=3$ \\
\hline
$\lambda_n$ & $4{\cdot}73004074$ & $10{\cdot}9956079$ & $17{\cdot}2787596$ \\
$\mu_n$ & $7{\cdot}8532046$ & $14{\cdot}1371655$ & $20{\cdot}4203522$ \\
$\cosh\frac{1}{2}\lambda_n$ & & & \\
$\cos\frac{1}{2}\lambda_n$ & & & \\
$\sinh\frac{1}{2}\mu_n$ & & & \\
$\sin\frac{1}{2}\mu_n$ & & & \\
\hline
\end{tabular}
\end{table}

\begin{equation}
\label{eq:A5-13}
\lambda_m \sim (2m - \tfrac{3}{2})\pi \quad \text{and} \quad \mu_m \sim (2m+\tfrac{1}{2})\pi \quad (m\to\infty)
\end{equation}
give the roots correct to ten decimals.

The functions $C_n$ and $S_n$ for $m = 1, 2, 3$, and 4 (together with their
first three derivatives) have been tabulated to nine decimals by Harris
and Reid for $x = 0\;(0{\cdot}005)\;0{\cdot}5$.

\paragraph{(a) Integrals involving $C_n(x)$ and $S_n(x)$}
In the application of the $C$- and $S$-functions, we have often to
evaluate matrix elements involving them. Simple integrals of the type
one generally encounters (as in \S\S~86--88) can be most readily evaluated
by raising by four the order of the derivative with which $C_n(x)$ (or
$S_n(x)$) occurs in the integrand (by making use of the differential equation~(\ref{eq:A5-1}) which they satisfy) and then following by a sequence of four integrations by parts. As illustrative of this method we shall consider two
typical examples.

(1) Consider
\begin{equation}
\label{eq:A5-14}
(S_n|C_m') = \int_{-\frac{1}{2}}^{+\frac{1}{2}} S_n(x) C_m'(x) \, dx.
\end{equation}

We first rewrite this integral as
\begin{equation}
\label{eq:A5-15}
\mu_n^4 (S_n|C_m') = \int_{-\frac{1}{2}}^{+\frac{1}{2}} S_n^{(4)}(x) C_m'(x) \, dx = (S_n''|C_m''').
\end{equation}

A first integration by parts now gives
\begin{equation}
\label{eq:A5-16}
\mu_n^4 (S_n|C_m') = S_n''({\tfrac{1}{2}}) C_m'''({\tfrac{1}{2}}) - \lambda_m^4 (S_n'|C_m').
\end{equation}

By a further sequence of integrations by parts, we find
\begin{equation}
\label{eq:A5-17}
(S_n'|C_m') = -(S_n''|C_m) = +(S_n|C_m'') = -(S_n|C_m'').
\end{equation}

Using this last result in equation~(\ref{eq:A5-16}), we obtain
\begin{equation}
\label{eq:A5-18}
(S_n|C_m') = \frac{2 S_n''(\frac{1}{2}) C_m'''(\frac{1}{2})}{\mu_n^4 + \lambda_m^4}.
\end{equation}

(2) Consider the matrix element
\begin{equation}
\label{eq:A5-19}
(S_n|x|C_m) = \int_{-\frac{1}{2}}^{+\frac{1}{2}} S_n \, x \, C_m' \, dx.
\end{equation}

In the same way as in the preceding example, we find
\begin{equation}
\label{eq:A5-20}
\lambda_m^4 (S_n|x|C_m') = (S_n|x|C_m^{(4)}) = -(S_n|C_m''') -(S_n'|C_m'') + (S_n''|C_m') + (S_n'''|C_m).
\end{equation}

We thus find (making further use of the relations included in~(\ref{eq:A5-17}))
\begin{equation}
\label{eq:A5-21}
(S_n|x|C_m') = \frac{-4 S_n''(\frac{1}{2}) C_m'''(\frac{1}{2})}{\left(\mu_n^4 + \lambda_m^4\right)^2} \cdot \lambda_m^4 S_n''(\frac{1}{2}) C_m'''(\frac{1}{2}).
\end{equation}

\section{Functions suitable for problems with cylindrical and spherical boundaries}
\label{sec:134}

For problems with cylindrical and spherical boundaries, the required
functions can be derived as the proper solutions of the characteristic
value problem defined by the equation
\begin{equation}
\label{eq:A5-22}
\mathscr{D}^4 y = \left(\frac{d^2}{dr^2} + \frac{1}{r}\frac{d}{dr} - \frac{\nu^2}{r^2}\right)^2 y = \alpha y,
\end{equation}
and the boundary conditions
\begin{equation}
\label{eq:A5-23}
y = 0 \quad \text{and} \quad dy/dr = 0 \quad \text{for} \quad r = 1 \quad \text{and} \quad r = \eta \;(< 1).
\end{equation}

First we may observe that if $f$ is any continuous and bounded function in the interval $(\eta, 1)$ and $y$ satisfies the boundary conditions~(\ref{eq:A5-23}),
then
\begin{equation}
\label{eq:A5-24}
\int_\eta^1 r \, y \, \mathscr{D}^4 f \, dr = \int_\eta^1 f \left(\frac{d^2}{dr^2} + \frac{1}{r}\frac{d}{dr} - \frac{\nu^2}{r^2}\right)^2 y \, dr
= -\int_\eta^1 f \, r \, \mathscr{D}^2 y \, dr.
\end{equation}

If $y_n$ and $y_m$ are proper solutions belonging to two different characteristic values $\alpha_n$ and $\alpha_m$, then it follows from the equation governing
$y_n$ that
\begin{equation}
\label{eq:A5-25}
\alpha_k \int_\eta^1 r \, y_n \, y_m \, dr = \int_\eta^1 r \, y_m \, \mathscr{D}^4 y_n \, dr.
\end{equation}

Now making use of the result expressed in~(\ref{eq:A5-24}), we obtain
\begin{equation}
\label{eq:A5-26}
\alpha_n \int_\eta^1 r \, y_n \, y_m \, dr = \int_\eta^1 r (\mathscr{D}^2 y_n)(\mathscr{D}^2 y_m) \, dr.
\end{equation}

From the symmetry of the integral on the right-hand side of equation~(\ref{eq:A5-26}) in $n$ and $m$ we conclude that
\begin{equation}
\label{eq:A5-27}
(\alpha_n - \alpha_m) \int_\eta^1 r \, y_n \, y_m \, dr = 0.
\end{equation}

Hence
\begin{equation}
\label{eq:A5-28}
\int_\eta^1 r \, y_n \, y_m \, dr = 0 \quad \text{if} \quad n \neq m.
\end{equation}

The proper solutions $y_n$, therefore, form an orthogonal set in the interval
$(1, \eta)$ with the weight function $r$.

Returning to equation~(\ref{eq:A5-22}), we may write its general solution as a
linear combination of Bessel functions of the different kinds and
arguments of order $\nu$; thus
\begin{equation}
\label{eq:A5-29}
y = A J_\nu(\alpha r) + B Y_\nu(\alpha r) + C I_\nu(\alpha r) + D K_\nu(\alpha r),
\end{equation}
where $A$, $B$, $C$, and $D$ are arbitrary constants. (In writing the solution
in the form~(\ref{eq:A5-29}) we have assumed that $\nu$ is an integer; if $\nu$ should not
be an integer we have only to write $J_{-\nu}$ in place of $Y_\nu$.)

From the recurrence relations satisfied by the Bessel functions, it
follows that
\begin{equation}
\label{eq:A5-30}
\frac{d}{dr}(ry) = A\alpha J_{\nu-1}(\alpha r) + B\alpha Y_{\nu-1}(\alpha r) + C\alpha I_{\nu-1}(\alpha r) - D\alpha K_{\nu-1}(\alpha r).
\end{equation}

The application of the boundary conditions~(\ref{eq:A5-23}) to the solution~(\ref{eq:A5-29}), leads to the characteristic equation
\begin{equation}
\label{eq:A5-31}
\begin{vmatrix}
J_\nu(\alpha) & Y_\nu(\alpha) & I_\nu(\alpha) & K_\nu(\alpha) \\
J_\nu(\alpha\eta) & Y_\nu(\alpha\eta) & I_\nu(\alpha\eta) & K_\nu(\alpha\eta) \\
J_{\nu-1}(\alpha) & Y_{\nu-1}(\alpha) & I_{\nu-1}(\alpha) & -K_{\nu-1}(\alpha) \\
J_{\nu-1}(\alpha\eta) & Y_{\nu-1}(\alpha\eta) & I_{\nu-1}(\alpha\eta) & -K_{\nu-1}(\alpha\eta)
\end{vmatrix} = 0.
\end{equation}

If $\alpha_m$ is a characteristic root and $(1, B_m, C_m, D_m)$ is the vector which is
annihilated by the matrix in~(\ref{eq:A5-31}) (when $\alpha = \alpha_m$), then
\begin{equation}
\label{eq:A5-32}
\mathscr{V}_{m}(r) = J_\nu(\alpha_m r) + B_m Y_\nu(\alpha_m r) + C_m I_\nu(\alpha_m r) + D_m K_\nu(\alpha_m r)
\end{equation}
is a proper solution belonging to $\alpha_m$.

\paragraph{(a) The normalization integral}

Let
\begin{equation}
\label{eq:A5-33}
u_{\nu m}(r) = J_\nu(\alpha_m r) + B_m Y_\nu(\alpha_m r)
\end{equation}
and
\begin{equation}
\label{eq:A5-34}
v_{\nu m}(r) = C_m I_\nu(\alpha_m r) + D_m K_\nu(\alpha_m r),
\end{equation}
so that
\begin{equation}
\label{eq:A5-35}
\mathscr{V}_{\nu m}(r) = u_{\nu m}(r) + v_{\nu m}(r).
\end{equation}

In virtue of the boundary conditions satisfied by $\mathscr{V}_{\nu m}$,
\begin{equation}
\label{eq:A5-36}
u_{\nu m}(1) = -v_{\nu m}(1); \quad u_{\nu m}'(1) = -v_{\nu m}'(1);
\end{equation}
\begin{equation}
\label{eq:A5-36a}
u_{\nu m}(\eta) = -v_{\nu m}(\eta); \quad u_{\nu m}'(\eta) = -v_{\nu m}'(\eta);
\end{equation}
and in particular,
\begin{equation}
\label{eq:A5-36b}
u_{\nu m} v_{\nu m}' - u_{\nu m}' v_{\nu m} = 0 \quad \text{at} \quad r = 1 \quad \text{and} \quad \eta.
\end{equation}

The functions $u_{\nu m}$ and $v_{\nu m}$, each being a linear combination of Bessel
functions with like (i.e.\@ real or imaginary) arguments, they satisfy the
differential equations
\begin{equation}
\label{eq:A5-37}
\frac{1}{r}\frac{d}{dr}\left(r\frac{du_{\nu m}}{dr}\right) - \frac{\nu^2}{r^2}u_{\nu m} = -\alpha_m^2 u_{\nu m}
\end{equation}
and
\begin{equation}
\label{eq:A5-38}
\frac{1}{r}\frac{d}{dr}\left(r\frac{dv_{\nu m}}{dr}\right) - \frac{\nu^2}{r^2}v_{\nu m} = +\alpha_m^2 v_{\nu m}.
\end{equation}

Multiplying equations~(\ref{eq:A5-37}) and~(\ref{eq:A5-38}) by $r v_{\nu m}$ and $r u_{\nu m}$, respectively,
and then subtracting one from another, we find
\begin{equation}
\label{eq:A5-39}
2\alpha_m^2 r \, u_{\nu m} \, v_{\nu m} = \frac{d}{dr}\left(r \, u_{\nu m} \frac{dv_{\nu m}}{dr} - r \, v_{\nu m} \frac{du_{\nu m}}{dr}\right),
\end{equation}
and integrating this equation over the range of $r$, we obtain
\begin{equation}
\label{eq:A5-40}
2\alpha_m^2 \int_\eta^1 r \, u_{\nu m} \, v_{\nu m} \, dr = \bigl[r(u_{\nu m} v_{\nu m}' - v_{\nu m} u_{\nu m}')\bigr]_\eta^1.
\end{equation}

The quantity on the right-hand side of equation~(\ref{eq:A5-40}) vanishes by virtue
of equation~(\ref{eq:A5-36b}). Hence,
\begin{equation}
\label{eq:A5-41}
\int_\eta^1 r \, u_{\nu m}^2 \, dr = 0.
\end{equation}

Making use of this property, we can write
\begin{equation}
\label{eq:A5-42}
\int_\eta^1 r \, \mathscr{V}_{\nu m}^2 \, dr = \int_\eta^1 r(u_{\nu m}^2 + v_{\nu m}^2) \, dr.
\end{equation}

To evaluate the integrals on the right-hand side, consider, first, the
differential equation satisfied by $u_{\nu m}$; after multiplication by $2r u_{\nu m}'$, the
equation can be written in the form
\begin{equation}
\label{eq:A5-43}
\frac{d}{dr}(r u_{\nu m}')^2 = -\alpha_m^2 r \frac{d r^2}{dr} u_{\nu m}^2 + \frac{\nu^2}{r}\frac{du_{\nu m}^2}{dr}.
\end{equation}

On integrating this equation over the range of $r$, we obtain
\begin{equation}
\label{eq:A5-44}
\bigl[r^2(u_{\nu m}')^2\bigr]_\eta^1 = -\alpha_m^2 \bigl[r^2 u_{\nu m}^2\bigr]_\eta^1 + 2\alpha_m^2 \int_\eta^1 r \, u_{\nu m}^2 \, dr + \bigl[\nu^2 u_{\nu m}^2\bigr]_\eta^1.
\end{equation}

We thus have
\begin{equation}
\label{eq:A5-45}
2\alpha_m^2 \int_\eta^1 r \, u_{\nu m}^2 \, dr = -r^2(u_{\nu m}')^2 + (\alpha_m^2 r^2 - \nu^2) u_{\nu m}^2 \Big|_\eta^1.
\end{equation}

From the equation satisfied by $v_{\nu m}$, we similarly find
\begin{equation}
\label{eq:A5-46}
2\alpha_m^2 \int_\eta^1 r \, v_{\nu m}^2 \, dr = \bigl[r^2(v_{\nu m}')^2 + (\alpha_m^2 r^2 - \nu^2) v_{\nu m}^2\bigr]_\eta^1.
\end{equation}

Now adding equations~(\ref{eq:A5-45}) and~(\ref{eq:A5-46}), we have
\begin{equation}
\label{eq:A5-47}
2\alpha_m^2 \int_\eta^1 r (u_{\nu m}^2 + v_{\nu m}^2) \, dr = \bigl[r^2\{(v_{\nu m}')^2 - (u_{\nu m}')^2\} +
(\alpha_m^2 r^2 - \nu^2)(v_{\nu m}^2 + u_{\nu m}^2)\bigr]_\eta^1.
\end{equation}

In view of the relations given in~(\ref{eq:A5-36}), equation~(\ref{eq:A5-47}) reduces to
\begin{equation}
\label{eq:A5-48}
\int_\eta^1 r (u_{\nu m}^2 + v_{\nu m}^2) \, dr = \bigl[r^2 u_{\nu m}^2\bigr]_\eta^1 = [r^2 v_{\nu m}^2]_\eta^1.
\end{equation}

Returning to equation~(\ref{eq:A5-42}) and remembering that the functions $\mathscr{V}_{\nu m}$
and $\mathscr{V}_{\nu n}$ for $m \neq n$ are orthogonal (in the interval $(1,\eta)$ and with the
weight function $r$), we can write
\begin{equation}
\label{eq:A5-49}
\int_\eta^1 r \, \mathscr{V}_{\nu m}^2 \, dr = N_{\nu m},
\end{equation}
where
\begin{equation}
\label{eq:A5-50}
N_{\nu m} = v_{\nu m}^2(1) - \eta^2 v_{\nu m}^2(\eta) = u_{\nu m}^2(1) - \eta^2 u_{\nu m}^2(\eta).
\end{equation}

\paragraph{(b) The cylinder functions of order 1}

Axisymmetric viscous flows between rotating cylinders can be expanded in terms of the functions of order~1; such an expansion was
used for example in Chapter~VII, \S~73. In view of the special importance
of the functions of order~1, we list in Table~LXIX a few of the characteristic roots and certain related constants for three different values
of~$\eta$.

For the particular case $\nu = \frac{1}{2}$, the solutions belonging to the first
three roots have been tabulated by Chandrasekhar and Elbert.

\begin{table}[htbp]
\centering
\caption{The characteristic roots and related constants for $\nu = 1$}
\label{tab:LXIX}
\begin{tabular}{l|ccc|ccc|c}
\hline
 & \multicolumn{3}{c|}{$\eta = \frac{1}{4}$} & \multicolumn{3}{c|}{$\eta = \frac{1}{2}$} & $\eta = \frac{3}{4}$ \\
 & $m=1$ & $m=2$ & $m=3$ & $m=1$ & $m=2$ & $m=3$ & $m=1$ \\
\hline
$\alpha_m$         & $+6{\cdot}390336$ & $+10{\cdot}55\hat{1}651$ & $+14{\cdot}729000$ & $+9{\cdot}49896$ & $+15{\cdot}739746$ & $+22{\cdot}01782\hat{1}$ & $+18{\cdot}93440$ \\
$B_m$              & $-5{\cdot}870144$ & $-0{\cdot}4281099$ & $+0{\cdot}7125113$ & $-9{\cdot}076420$ & $-17{\cdot}3573$ & $-22{\cdot}4667$ & $-10{\cdot}2780$ \\
$C_m$              & $-1{\cdot}497055\times10^{-2}$ & $+4{\cdot}167599\times10^{-5}$ & $-7{\cdot}144507\times10^{-7}$ & $+1{\cdot}006354\times10^{-3}$ & $+3{\cdot}67581\times10^{-6}$ & $+8{\cdot}86448\times10^{-9}$ & $-8{\cdot}91136\times10^{-3}$ \\
$D_m$              & $-10{\cdot}77988$ & $-5{\cdot}986757$ & $-19{\cdot}90159$ & $+4{\cdot}396374\times10^{2}$ & $-1{\cdot}95381\times10^{4}$ & $+5{\cdot}91172\times10^{5}$ & $-6{\cdot}65593\times10^{8}$ \\
$u_m(1) = -v_m(1)$ & $+1{\cdot}330279$ & $-0{\cdot}1884864$ & $+0{\cdot}1803302$ & $-1{\cdot}681248$ & $-2{\cdot}46940$ & $-2{\cdot}70290$ & $+1{\cdot}34981$ \\
$u_m(\eta) = -v_m(\eta)$ & $+2{\cdot}619034$ & $+0{\cdot}3728952$ & $+0{\cdot}3581128$ & $-2{\cdot}371094$ & $+3{\cdot}48706$ & $-3{\cdot}81917$ & $+1{\cdot}55828$ \\
$u'_m(1) = -v'_m(1)$ & $+7{\cdot}807580$ & $-1{\cdot}903352$ & $+2{\cdot}570976$ & $-14{\cdot}93868$ & $-37{\cdot}7250$ & $-58{\cdot}2068$ & $+24{\cdot}4679$ \\
$u'_m(\eta) = -v'_m(\eta)$ & $-23{\cdot}29673$ & $-4{\cdot}83996$ & $-6{\cdot}10685$ & $+24{\cdot}82964$ & $-58{\cdot}7097$ & $+88{\cdot}1452$ & $-30{\cdot}0835$ \\
$N_{1;m}$          & $+1{\cdot}340934$ & $+0{\cdot}02683645$ & $+0{\cdot}02450369$ & $+1{\cdot}42107$ & $+3{\cdot}05805$ & $+3{\cdot}65915$ & $+0{\cdot}456104$ \\
\hline
\end{tabular}
\end{table}

\paragraph{(c) The spherical functions of half-odd integral orders}

The functions of half-odd integral orders are useful in problems of
viscous flow in spheres and in spherical shells. For flows in a sphere,
the requirement that the solutions have no singularity at the origin
implies that in the solution~(\ref{eq:A5-29}) we cannot include the terms in $Y_\nu$
(or $J_{-\nu}$) and $K_\nu$. Therefore, in this case we write
\begin{equation}
\label{eq:A5-51}
\mathscr{V}_{\nu,l}(r) = \frac{J_\nu(\alpha_l r)}{J_\nu(\alpha_l)} - \frac{I_\nu(\alpha_l r)}{I_\nu(\alpha_l)}.
\end{equation}

Defined in this manner, the function $\mathscr{V}_{\nu,l}(r)$ clearly vanishes at $r = 1$;
and the condition that its derivative also vanishes at $r = 1$ leads to
the characteristic equation
\begin{equation}
\label{eq:A5-52}
J_\nu'(\alpha_l) I_\nu(\alpha_l) - I_\nu'(\alpha_l) J_\nu(\alpha_l) = 0.
\end{equation}

The first three roots of this equation for $l = 1, 3$, and 5 are listed in
Table~LXX.

The roots of the equation
\begin{equation}
\label{eq:A5-53}
J_\nu(\alpha_l) I_\nu'(\alpha_l) - I_\nu(\alpha_l) J_\nu'(\alpha_l) = \alpha_l I_\nu(\alpha_l) J_\nu(\alpha_l)
\end{equation}
needed in a related connexion (\S~62) are also included in Table~LXX.

\begin{table}[htbp]
\centering
\caption{The roots of equations~(\ref{eq:A5-52}) and~(\ref{eq:A5-53})}
\label{tab:LXX}
\begin{tabular}{l|ccc|ccc}
\hline
 & \multicolumn{3}{c|}{Roots of equation~(\ref{eq:A5-52})} & \multicolumn{3}{c}{Roots of equation~(\ref{eq:A5-53})} \\
 & $l = 1$ & $l = 3$ & $l = 5$ & $l = 1$ & $l = 3$ & $l = 5$ \\
\hline
$\frac{3}{2}$ & $5{\cdot}763$ & $12{\cdot}065$ & $18{\cdot}346$ & $6{\cdot}913$ & $13{\cdot}363$ & $19{\cdot}653$ \\
$\frac{5}{2}$ & $7{\cdot}458$ & $13{\cdot}680$ & $19{\cdot}943$ & $8{\cdot}574$ & $14{\cdot}870$ & $21{\cdot}112$ \\
$\frac{7}{2}$ & $9{\cdot}108$ & $15{\cdot}385$ & $21{\cdot}620$ & $10{\cdot}209$ & $16{\cdot}410$ & \\
\hline
\end{tabular}
\end{table}

\section*{BIBLIOGRAPHICAL NOTES}

The usefulness of the functions described in this appendix for the solution of
hydrodynamical problems was pointed out by:

\begin{enumerate}
\item S. \textsc{Chandrasekhar} and W.~H. \textsc{Reid}, `On the expansion of functions which
satisfy four boundary conditions', \textit{Proc. Nat. Acad. Sci.} \textbf{43}, 521--7 (1957).
\end{enumerate}

\noindent\S~\textbf{133}. The functions which we have denoted by $C_n$ and $S_n$ occur in the theory of
vibrating beams and in that connexion they have been considered by Rayleigh:

\begin{enumerate}
\setcounter{enumi}{1}
\item Lord \textsc{Rayleigh}, \textit{Theory of Sound}, p.~278, Dover Publications, New York,
1945.
\end{enumerate}

\noindent In the hydrodynamical context these functions have been studied most extensively by:

\begin{enumerate}
\setcounter{enumi}{2}
\item W.~H. \textsc{Reid}, `On the stability of viscous flow in a curved channel', \textit{Proc.
Roy. Soc. (London)} A, \textbf{244}, 186--98 (1958).

\item D.~L. \textsc{Harris} and W.~H. \textsc{Reid}, `On orthogonal functions which satisfy
four boundary conditions. I. Tables for use in Fourier-type expansions',
\textit{Astrophys. J. Suppl. Ser.} \textbf{3}, 429--47 (1958).

\item W.~H. \textsc{Reid} and D.~L. \textsc{Harris}, `On orthogonal functions which satisfy
four boundary conditions.  II. Integrals for use with Fourier-type expansions', \textit{Astrophys. J. Suppl. Ser.} \textbf{3}, 448--52 (1958).
\end{enumerate}

\noindent In reference~4 Harris and Reid have tabulated the functions $C_m$ and $S_m$ and their
first three derivatives for $m = 1, \ldots, 4$ to 9 decimals for $x = 0\;(0{\cdot}005)\;0{\cdot}5$. Less
extensive tables have been published earlier by:

\begin{enumerate}
\setcounter{enumi}{5}
\item R.~E.~D. \textsc{Bishop} and D.~C. \textsc{Johnson}, \textit{Vibration Analysis Tables}, Cambridge
University Press, New York, 1956.
\end{enumerate}

\noindent In reference~5 Reid and Harris have provided a very useful compilation of
integrals involving the $C$- and the $S$-functions.

\noindent\S~\textbf{134}. The functions denoted by $\mathscr{V}_{\nu;m}$ seem to have been considered for the first
time in these connexions in reference~1. See also:

\begin{enumerate}
\setcounter{enumi}{6}
\item S. \textsc{Chandrasekhar}, `The stability of viscous flow between rotating
cylinders', \textit{Proc. Roy. Soc. (London)} A, \textbf{246}, 301--11 (1958).
\end{enumerate}

\noindent The numerical results included in Table~LXIX are taken from:

\begin{enumerate}
\setcounter{enumi}{7}
\item \textemdash\textemdash{} and Donna D. \textsc{Elbert}, `On orthogonal functions which satisfy four
boundary conditions. III. Tables for use in Fourier--Bessel-type expansions', \textit{Astrophys. J. Suppl. Ser.} \textbf{3}, 453--8 (1958).
\end{enumerate}

\noindent The functions $\mathscr{V}_{1;m}$ for $m = 1, 2$, and 3 and $\eta = \frac{1}{2}$ are tabulated in reference~8.

\noindent The characteristic roots listed in Table~LXX are taken from:

\begin{enumerate}
\setcounter{enumi}{8}
\item F.~E. \textsc{Bisshopp}, `On the thermal instability of a rotating fluid sphere',
\textit{Phil. Mag.} Ser. 8, \textbf{3}, 1342--60 (1958).
\end{enumerate}
